\documentclass[a4paper]{article}     
\usepackage{amsfonts}
\usepackage{amsmath}
\usepackage{amssymb}
\usepackage{graphicx}

\begin{document}
	
	\noindent \large Auteur: Victor Pena \\
	\noindent \large Studierichting: Informatica\\
	\noindent \large Oefeningenreeks 4A nr. 46.2\\
	
	\medskip
	
	\normalsize 
	
	$\displaystyle\int\dfrac{-3x+3}{2x^4+5x^2+2}dx \overset{u=x^2}{\Rightarrow} u = \dfrac{-5 \pm \sqrt{5^2-4(2)(2)}}{2(2)} = \left\{-2,-\dfrac{1}{2}\right\} \Rightarrow \dfrac{1}{2}\displaystyle\int \dfrac{-3x+3}{\left(x^2+2\right)\left(x^2+\dfrac{1}{2}\right) }dx = $\\
	
	$\dfrac{1}{2} \displaystyle\int\dfrac{Ax+B}{(x^2+2)}dx + \dfrac{1}{2}\displaystyle\int \dfrac{Cx+D}{\left(x^2+\dfrac{1}{2}\right)}dx $ \\
	
	$$\Rightarrow \begin{array}{ccc}
		A\left(i\sqrt{2}\right) + B=&\left.\dfrac{-3x+3}{\left(x+\dfrac{1}{2}\right)}\right|_{x=i\sqrt{2}}&=-2+i(2\sqrt{2}) \Rightarrow A = 2, B = -2 \\
		C\left(\dfrac{i}{\sqrt{2}}\right) + D=&\left.\dfrac{-3x+3}{(x^2+2)}\right|_{x=\tfrac{i}{\sqrt{2}}}&=2-i\sqrt{2} \Rightarrow C =-2, D = 2 \\
	\end{array}$$ \\
	
	$\Rightarrow \dfrac{1}{2} \displaystyle\int \dfrac{2x-2}{x^2+2}dx + \dfrac{1}{2} \displaystyle\int \dfrac{-2x+2}{x^2+\dfrac{1}{2}}dx = $\\
	
	$ \dfrac{1}{2} \displaystyle\int \dfrac{d(x^2+2)}{x^2+2} - \dfrac{1}{\sqrt{2}} \displaystyle\int\dfrac{d\left(\dfrac{x}{\sqrt{2}}\right)}{\left(\dfrac{x}{\sqrt{2}}\right)^2 + 1} - \dfrac{1}{2} \displaystyle\int \dfrac{d\left(x^2+\dfrac{1}{2}\right)}{x^2+\dfrac{1}{2}} + \sqrt{2} \displaystyle\int \dfrac{d\left(x\sqrt{2}\right)}{(x\sqrt{2})^2+1}$ \\
	
	$= \dfrac{1}{2} \ln|x^2+2| - \dfrac{1}{\sqrt{2}} \operatorname{Bgtan}\left(\dfrac{x}{\sqrt{2}}\right) - \dfrac{1}{2} \ln\left|x^2+\dfrac{1}{2}\right| + \sqrt{2}\operatorname{Bgtan}\left(x\sqrt{2}\right) + k$
	
	
	\begin{thebibliography}{999}
		%\bibitem{a} Referentie
	\end{thebibliography}
\end{document} 
